% !TeX program = xelatex
% ============================================================================
%  第 7 课　群里的小世界：阶与子群　—— 教师版讲义
%  与 02-课件/第07课-阶与子群/第07课-阶与子群.tex 同步
% ============================================================================

\jhlesson{7}{群里的小世界：阶与子群}{时长 45 分钟\quad·\quad 教具：正三角形纸片\quad·\quad 本课不发单页}

\begin{mubiao}
  \begin{itemize}
    \item 知道元素的\textbf{阶}就是「这个动作重复几次回到原位」，能从两行式里读出来。
    \item 能在 \(S_3\) 里圈出能自成一体的小群（子群）。
    \item 记住检查子群的办法：看它会不会「跑出去」。
  \end{itemize}
\end{mubiao}

\noindent{\small 本课节奏：0—6 回顾｜6—20 元素的阶｜20—40 子群｜40—45 收尾。全程至少 4 次「用手指回答」。}

\section*{一、回顾（0—6 分钟）}

把上节课乘法表的两行再念一遍（\(r\) 行、\(f\) 行），指出它们都是同样 \(6\) 个动作的一次重排。

一句话过渡：今天换一个角度——不看一整行，只看\emph{一个动作自己反复做}会怎么样。

\section*{二、元素的阶（6—20 分钟）}

\begin{definition}[元素的阶]
  一个动作 \(g\) 连续做 \(n\) 次，\textbf{第一次}回到「不动」，这个 \(n\)
  就叫做 \(g\) 的\textbf{阶}。
\end{definition}

\begin{caozuo}
  拿纸片当场连做给学生看，边做边数：
  \begin{itemize}
    \item \(e\)（不动）：做 \(1\) 次就还是 \(e\)，阶是 \(1\)；
    \item \(r\)（转 \(120^\circ\)）：\(r\to r^2\to e\)，做 \(3\) 次回来，阶是 \(3\)；
    \item \(f\)（沿过 \(A\) 的轴翻）：\(f\to e\)，做 \(2\) 次回来，阶是 \(2\)。
  \end{itemize}
  再补一句：三个翻折的阶都是 \(2\)，两个转动的阶都是 \(3\)。
\end{caozuo}

白话解释一遍：阶 \(1\) 是「做一次就回到原位」，阶 \(2\) 是「做两次」，阶 \(3\) 是「做三次」。
「阶」就是「这个动作要重复几次，才转完一圈回到自己」。

\subsection*{顺便看一眼循环记法}

转 \(120^\circ\) 的两行式可以缩写成一个圈：

\[
  \begin{pmatrix}A&B&C\\B&C&A\end{pmatrix}
  \quad\rightsquigarrow\quad
  (A\ B\ C)
\]

\begin{zhu}
  \textbf{只提一句就好。}说清楚两件事：读法是「\(A\) 去 \(B\)，\(B\) 去 \(C\)，\(C\) 回 \(A\)」；
  圈里有几个字母，阶就是几。不要求学生会写会用，了解即可。
\end{zhu}

\begin{caozuo}
  手指回答：\(r^2\) 的阶是多少？（\(3\)）\(rf\) 的阶是多少？（\(2\)）\((A\ B)\) 的阶是多少？（\(2\)）
\end{caozuo}

\section*{三、子群（20—40 分钟）}

\begin{definition}[子群]
  设 \(H\) 是群 \(G\) 的一部分。如果 \(H\) 按同一个运算\textbf{自己也构成一个群}，
  就说 \(H\) 是 \(G\) 的一个\textbf{子群}。
\end{definition}

\begin{zhu}
  检查的办法只讲一条，也只用一条：\textbf{会不会跑出去}。
  把这张小表填满，看结果是不是全都落在圈定的范围里。
\end{zhu}

\begin{caozuo}
  一步一步带着全班圈，每圈一个就填一张小表：
  \begin{enumerate}
    \item \(H=\{e,r,r^2\}\)（三个转动）：填表后结果全在里面，\textbf{是}子群；
    \item \(H=\{e,f\}\)：\(f\) 连做两次回到 \(e\)，全在里面，\textbf{是}子群；
    \item \(H=\{e,r,f\}\)：\(r\cdot r=r^2\) 不在里面，\(r\cdot f=rf\) 也不在里面，
      \textbf{不是}子群。
  \end{enumerate}
  再让全班自己试 \(\{e,rf\}\)、\(\{e,r^2f\}\)，都会发现是子群。
\end{caozuo}

把 \(S_3\) 的子群列成一张表：

\begin{center}
\begin{tabular}{@{}llc@{}}
  \toprule
  子群 & 是谁 & 大小 \\
  \midrule
  \(\{e\}\) & 只有不动 & \(1\) \\
  \(\{e,\ r,\ r^2\}\) & 三个转动 & \(3\) \\
  \(\{e,\ f\}\) & 不动加一条轴的翻折 & \(2\) \\
  \(\{e,\ rf\}\) & 同上 & \(2\) \\
  \(\{e,\ r^2f\}\) & 同上 & \(2\) \\
  \(S_3\) & 全部六个 & \(6\) \\
  \bottomrule
\end{tabular}
\end{center}

\begin{zhu}
  这里先埋一个观察：小群的大小 \(1,2,3,6\)，每一个都能整除大群的大小 \(6\)。
  今天只要学生点头，不要求解释。第 10 课会在模 \(12\) 加法里再看一遍同样的事。
\end{zhu}

\section*{四、收尾（40—45 分钟）}

板书两句话：元素的阶，是这个动作重复几次回到原位；
群里面还能装小群，检查的办法是看它会不会跑出去。

\begin{xiaojie}
  \begin{itemize}
    \item 元素的\textbf{阶}：\(e\) 是 \(1\)，三个翻折是 \(2\)，两个转动是 \(3\)。
    \item \textbf{子群}：\(S_3\) 里有 \(\{e\}\)、\(\{e,r,r^2\}\) 和三个 \(\{e,\text{翻折}\}\) 共 \(5\) 个真子群。
    \item 子群的大小 \(1,2,3\) 都能整除 \(S_3\) 的大小 \(6\)。
  \end{itemize}
\end{xiaojie}

\noindent\textbf{下节课}：换一张更大的纸——正方形的 \(8\) 个对称动作。